% =======================================================================
%  Esquema del banco de medicion coherente DP-QPSK
%  Reemplaza el PNG lab_setup_1.png por una version vectorial editable.
%
%  Cambios respecto del dibujo anterior:
%    - despues del tap 95/5 se agrego el atenuador variable manual
%    - despues del atenuador va un acoplador Newport SMF de 1550 nm, del
%      que se usa una sola entrada; una salida va al medidor de potencia y
%      la otra al front end coherente
%    - la relacion entre las dos salidas es la calibracion del banco,
%      P_rx = P_med + 26,32 dB, medida con -33,55 / -7,23 dBm
%    - el laser transmisor es un IDPhotonics CoBrite Tunable Laser DX
%    - el rotulo del acoplador va arriba porque abajo pasa la rama al medidor
%    - el rotulo del detector InGaAs paso a la izquierda por espacio
%
%  Compilar:  pdflatex lab_setup_1.tex
% =======================================================================
\documentclass[border=6pt]{standalone}
\usepackage[utf8]{inputenc}
\usepackage[T1]{fontenc}
\usepackage[spanish,es-noshorthands]{babel}
\usepackage{lmodern}
\usepackage{tikz}
\usetikzlibrary{arrows.meta, positioning, calc, fit, backgrounds}

\definecolor{cOpt}{RGB}{214,137,16}      % camino optico
\definecolor{cEle}{RGB}{26,49,77}        % camino electrico
\definecolor{cBorde}{RGB}{188,196,206}
\definecolor{cRelleno}{RGB}{246,248,251}
\definecolor{cPant}{RGB}{23,36,54}
\definecolor{cTrazo}{RGB}{26,49,77}
\definecolor{cAcento}{RGB}{56,129,176}

\begin{document}
\begin{tikzpicture}[
  x=0.92cm, y=0.92cm,
  font=\sffamily\small,
  caja/.style={draw=cBorde, fill=cRelleno, rounded corners=3pt,
               minimum width=2.05cm, minimum height=1.45cm, line width=0.8pt},
  rotulo/.style={font=\sffamily\small, align=center, text=cTrazo},
  opt/.style={draw=cOpt, line width=1.5pt, -{Stealth[length=3mm,width=2mm]}},
  ele/.style={draw=cEle, line width=1.2pt, -{Stealth[length=3mm,width=2mm]}},
  trig/.style={draw=cEle, line width=1.0pt, dash pattern=on 4pt off 2.5pt,
               -{Stealth[length=3mm,width=2mm]}},
  pant/.style={draw=cTrazo, fill=cPant, line width=0.7pt, rounded corners=1pt},
]

%% ---------------------------------------------------------------- nodos
\node[caja] (ppg)  at ( 0.0,  0.0) {};
\node[caja] (drv)  at ( 3.0,  0.0) {};
\node[caja] (mod)  at ( 6.0,  0.0) {};
\node[caja] (tap)  at ( 9.0,  0.0) {};
\node[caja] (atn)  at (11.9,  0.0) {};
\node[caja] (acop) at (14.9,  0.0) {};
\node[caja, minimum width=2.6cm] (fe) at (18.3, 0.0) {};
\node[caja] (adc)  at (21.5,  0.0) {};
\node[caja] (pc)   at (24.5,  0.0) {};

\node[caja] (ltx)  at ( 6.0,  3.3) {};
\node[caja] (llo)  at (18.3,  3.3) {};

\node[caja] (det)  at ( 9.0, -2.7) {};
\node[caja] (osc)  at ( 9.0, -5.3) {};
\node[caja] (dc)   at ( 5.2, -5.3) {};
\node[caja] (med)  at (14.9, -2.7) {};

%% ------------------------------------------------------------- iconos
% generador de patrones: pantalla con onda cuadrada y cuatro salidas
\draw[pant] ($(ppg)+(-0.62,-0.18)$) rectangle ($(ppg)+(0.12,0.42)$);
\draw[cAcento, line width=0.8pt]
  ($(ppg)+(-0.54,0.02)$) -- ++(0.14,0) -- ++(0,0.26) -- ++(0.16,0)
  -- ++(0,-0.26) -- ++(0.14,0) -- ++(0,0.26) -- ++(0.16,0) -- ++(0,-0.26) -- ++(0.1,0);
\foreach \k in {0,1,2,3}{
  \draw[cTrazo, line width=0.6pt, fill=white]
    ($(ppg)+(0.62,0.42-0.24*\k)$) circle (0.07);
}
\draw[cTrazo, line width=0.6pt] ($(ppg)+(-0.5,-0.42)$) -- ++(0.55,0);

% drivers: triangulo
\draw[cTrazo, line width=0.9pt, fill=white]
  ($(drv)+(-0.55,-0.6)$) -- ($(drv)+(-0.55,0.6)$) -- ($(drv)+(0.65,0)$) -- cycle;

% modulador: dos brazos Mach-Zehnder
\draw[cAcento, line width=0.9pt]
  ($(mod)+(-0.7,0)$) -- ++(0.22,0.42) -- ++(0.96,0) -- ++(0.22,-0.42)
  -- ++(-0.22,-0.42) -- ++(-0.96,0) -- cycle;
\draw[cTrazo, line width=2pt] ($(mod)+(-0.28,0.56)$) -- ++(0.56,0);
\draw[cTrazo, line width=2pt] ($(mod)+(-0.28,-0.56)$) -- ++(0.56,0);

% tap: acoplador
\draw[cTrazo, line width=0.8pt, fill=white, rounded corners=1.5pt]
  ($(tap)+(-0.22,-0.2)$) rectangle ($(tap)+(0.22,0.2)$);
\node[font=\sffamily\tiny, text=gray] at ($(tap)+(0.66,0.42)$) {95\,\%};
\node[font=\sffamily\tiny, text=gray] at ($(tap)+(0.5,-0.45)$) {5\,\%};

% front end: hibrido de 90 grados y balanceados
\draw[cAcento, line width=0.8pt, fill=white]
  ($(fe)+(-1.05,-0.55)$) rectangle ($(fe)+(-0.78,0.55)$);
\foreach \s in {1,-1}{
  \draw[cAcento, line width=0.7pt, fill=white]
    ($(fe)+(-0.35,0.42*\s-0.14)$) rectangle ($(fe)+(0.25,0.42*\s+0.2)$);
  \node[font=\sffamily\tiny, text=cAcento] at ($(fe)+(-0.05,0.42*\s+0.03)$) {$90^\circ$};
  \draw[cTrazo, line width=0.8pt, fill=cTrazo]
    ($(fe)+(0.5,0.42*\s+0.2)$) -- ($(fe)+(0.5,0.42*\s-0.14)$)
    -- ($(fe)+(0.82,0.42*\s+0.03)$) -- cycle;
  \draw[cOpt, line width=0.7pt] ($(fe)+(-0.78,0.42*\s+0.03)$) -- ($(fe)+(-0.35,0.42*\s+0.03)$);
  \draw[cTrazo, line width=0.7pt] ($(fe)+(0.82,0.42*\s+0.03)$) -- ($(fe)+(1.2,0.42*\s+0.03)$);
}

% laseres
\foreach \n in {ltx,llo}{
  \draw[cTrazo, line width=0.8pt, fill=white]
    ($(\n)+(-0.6,-0.4)$) rectangle ($(\n)+(-0.05,0.4)$);
  \fill[cPant] ($(\n)+(-0.48,-0.24)$) rectangle ($(\n)+(-0.17,0.24)$);
  \fill[cOpt, opacity=0.35] ($(\n)+(-0.05,0)$) -- ($(\n)+(0.7,0.3)$) -- ($(\n)+(0.7,-0.3)$) -- cycle;
  \draw[cOpt, line width=0.9pt] ($(\n)+(-0.05,0)$) -- ($(\n)+(0.7,0)$);
}

% detector InGaAs: fotodiodo
\draw[cTrazo, line width=0.9pt] ($(det)+(-0.45,0.45)$) -- ($(det)+(-0.45,-0.45)$);
\draw[cTrazo, line width=0.9pt, fill=cTrazo]
  ($(det)+(-0.45,0.34)$) -- ($(det)+(-0.45,-0.34)$) -- ($(det)+(0.12,0)$) -- cycle;
\draw[cTrazo, line width=0.9pt] ($(det)+(0.12,0.45)$) -- ($(det)+(0.12,-0.45)$);
\draw[cTrazo, line width=0.8pt] ($(det)+(0.12,0)$) -- ($(det)+(0.6,0)$);

% osciloscopios y PC: pantallas
\foreach \n/\dx in {osc/0, adc/0}{
  \draw[pant] ($(\n)+(-0.6,-0.42)$) rectangle ($(\n)+(0.45,0.42)$);
}
\draw[cAcento, line width=0.8pt]
  plot[domain=-0.5:0.35, samples=40, smooth]
  (\x+9.0, {0.18*sin(deg((\x+0.5)*9))-5.3});
\foreach \k in {0,1,2}{
  \draw[cTrazo, line width=0.6pt, fill=white] ($(osc)+(0.68,0.3-0.3*\k)$) circle (0.08);
  \draw[cTrazo, line width=0.6pt, fill=white] ($(adc)+(0.68,0.3-0.3*\k)$) circle (0.08);
  \draw[cTrazo, line width=0.6pt, fill=white] ($(adc)+(-0.78,0.3-0.3*\k)$) circle (0.08);
}
\foreach \k/\c in {0/cAcento, 1/cOpt}{
  \draw[\c, line width=0.7pt]
    plot[domain=-0.5:0.35, samples=40, smooth]
    (\x+21.5, {0.13*sin(deg((\x+0.5)*9+\k*90))+0.12-0.24*\k});
}

% PC
\draw[pant] ($(pc)+(-0.6,-0.24)$) rectangle ($(pc)+(0.6,0.5)$);
\foreach \a/\b in {-0.28/0.34, 0.28/0.34, -0.28/-0.06, 0.28/-0.06}{
  \fill[cAcento] ($(pc)+(\a,\b)$) circle (0.08);
}
\draw[cTrazo, line width=0.9pt] ($(pc)+(0,-0.24)$) -- ++(0,-0.2);
\draw[cTrazo, line width=0.9pt] ($(pc)+(-0.3,-0.44)$) -- ++(0.6,0);

% fuentes DC
\draw[pant] ($(dc)+(-0.6,0.06)$) rectangle ($(dc)+(0.12,0.46)$);
\foreach \k in {0,1,2}{
  \draw[cAcento, line width=1.1pt] ($(dc)+(-0.5+0.2*\k,0.26)$) -- ++(0.12,0);
}
\draw[cTrazo, line width=0.7pt, fill=white] ($(dc)+(0.42,0.26)$) circle (0.2);
\draw[cTrazo, line width=0.7pt] ($(dc)+(0.42,0.26)$) -- ++(0,0.15);
\foreach \k/\c in {0/red, 1/cTrazo, 2/red, 3/cTrazo}{
  \draw[\c, line width=0.7pt, fill=white] ($(dc)+(-0.48+0.32*\k,-0.4)$) circle (0.11);
}


% atenuador variable manual: pad con la flecha diagonal del ajuste
\draw[cTrazo, line width=0.9pt, fill=white]
  ($(atn)+(-0.45,-0.35)$) rectangle ($(atn)+(0.45,0.35)$);
\foreach \k in {0,1,2}{
  \draw[cTrazo, line width=0.7pt] ($(atn)+(-0.22+0.22*\k,-0.35)$) -- ++(0,0.7);
}
\draw[cAcento, line width=1.3pt, -{Stealth[length=3.4mm,width=2.4mm]}]
  ($(atn)+(-0.78,-0.66)$) -- ($(atn)+(0.78,0.66)$);

% acoplador: dos entradas, la de arriba sin usar, y dos salidas
\draw[cTrazo, line width=0.8pt, fill=white, rounded corners=1.5pt]
  ($(acop)+(-0.24,-0.3)$) rectangle ($(acop)+(0.28,0.3)$);
\draw[cOpt, line width=0.9pt] ($(acop)+(-0.62,-0.15)$) -- ($(acop)+(-0.24,-0.15)$);
\draw[cBorde, line width=0.9pt] ($(acop)+(-0.62,0.15)$) -- ($(acop)+(-0.24,0.15)$);
\draw[cTrazo, line width=0.7pt] ($(acop)+(-0.78,0.07)$) -- ++(0.16,0.16);
\draw[cTrazo, line width=0.7pt] ($(acop)+(-0.78,0.23)$) -- ++(0.16,-0.16);

% medidor de potencia optica: cabezal y display
\draw[cTrazo, line width=0.8pt, fill=white] ($(med)+(-0.72,0.1)$) circle (0.19);
\draw[cTrazo, line width=0.8pt] ($(med)+(-0.53,0.1)$) -- ($(med)+(-0.34,0.1)$);
\draw[pant] ($(med)+(-0.34,-0.12)$) rectangle ($(med)+(0.5,0.34)$);
\foreach \k in {0,1,2}{
  \draw[cAcento, line width=1.1pt] ($(med)+(-0.22+0.22*\k,0.11)$) -- ++(0.14,0);
}
\foreach \k in {0,1,2}{
  \draw[cTrazo, line width=0.6pt, fill=white] ($(med)+(-0.4+0.4*\k,-0.42)$) circle (0.09);
}

%% ------------------------------------------------------------- enlaces
\draw[ele] (ppg) -- (drv);
\draw[ele] (drv) -- (mod);
\draw[opt] (mod) -- (tap);
\draw[opt] (tap)  -- (atn);
\draw[opt] (atn)  -- (acop);
\draw[opt] (acop) -- (fe);
\draw[opt] (acop) -- (med);
\draw[ele] (fe)  -- (adc);
\draw[ele] (adc) -- (pc);

\draw[opt] (ltx) -- (mod);
\draw[opt] (llo) -- (fe);

% rama del 5 %: tap -> detector -> osciloscopio -> fuentes DC -> modulador
\draw[opt] (tap) -- (det);
\draw[ele] (det) -- (osc);
\draw[ele] (osc) -- (dc);
\draw[ele] (dc.north) -- ++(0,1.15) -| (mod.south);

% trigger del generador de patrones al ADC
\draw[trig] (ppg.west) -- ++(-1.0,0) -- ++(0,-7.9) -- (23.25,-7.9)
            -- (23.25,-0.40) -- ($(adc.east)+(0,-0.40)$);
\node[font=\sffamily\footnotesize, text=cEle, fill=white, inner sep=2pt]
  at (3.4,-7.9) {trigger de patr\'on};

%% ------------------------------------------------------------- rotulos
\node[rotulo, below=0.30 of ppg] {Generador de PRBS};
\node[rotulo, below=0.30 of drv] {Driver RF};
\node[rotulo, below=0.30 of mod] {Modulador DP-IQ};
\node[rotulo, below=0.30 of tap] {Tap 95/5};
\node[rotulo, below=0.30 of atn] {Atenuador\\variable};

\node[rotulo, below=0.30 of fe]  {Receptor coherente};
\node[rotulo, below=0.30 of adc] {Osciloscopio\\de tiempo real\\(4 canales)};
\node[rotulo, below=0.30 of pc]  {DSP offline};

\node[rotulo, above=0.30 of acop] {Acoplador};
\node[rotulo, above=0.30 of ltx] {L\'aser Tx\\Tunable Laser DX};
\node[rotulo, above=0.30 of llo] {L\'aser LO\\PPCL500};

\node[rotulo, left=0.30 of det] {Detector\\InGaAs};
\node[rotulo, below=0.30 of osc] {Osciloscopio\\de bias};
\node[rotulo, below=0.30 of dc]  {Fuentes DC bias};
\node[rotulo, below=0.30 of med] {Medidor de potencia};
\node[font=\sffamily\footnotesize, text=cTrazo, anchor=west, align=left]
  at ($(med)+(1.35,0.05)$) {$P_\mathrm{rx} = P_\mathrm{med} + 26{,}32$\,dB};

\end{tikzpicture}
\end{document}
